% Part: applied-modal-logics
% Chapter: epistemic-logic
% Section: introduction

\documentclass[../../../include/open-logic-section]{subfiles}

\begin{document}

\olfileid{aml}{el}{int}

\olsection{પરિચય}

મોડલ તર્ક \emph{મોડલ વિધાનો} અને તેમની વચ્ચેના તાર્કિક પરિણામોના
સંબંધોનો અભ્યાસ કરે છે; તે જ રીતે જ્ઞાનસંબંધી તર્ક
\emph{જ્ઞાનસંબંધી વિધાનો} અને તેમની વચ્ચેના તાર્કિક પરિણામોના
સંબંધોનો અભ્યાસ કરે છે. મોડલ સંક્રિયકોને શક્યતા અને અનિવાર્યતા
દર્શાવતા ગણવાને બદલે, અહીં એકસ્થાની સંયોજકોનું અર્થઘટન જ્ઞાન અને
માન્યતાના નિદર્શન માટે જ્ઞાનસંબંધી અથવા માન્યતાસંબંધી રીતે થાય છે.
દાખલા તરીકે, આપણે નીચેના જેવા દાવા વ્યક્ત કરવા માગી શકીએ:

\begin{enumerate}
\item રિચાર્ડ જાણે છે કે કૅલ્ગરી આલ્બર્ટામાં છે.
\item ઓડ્રીને લાગે છે કે સોફા પર કૂતરો હોઈ શકે છે.
\item રિચાર્ડ જાણે છે કે ઓડ્રી જાણે છે કે તેનો વર્ગ મંગળવારે હોય છે.
\item દરેક જણ જાણે છે કે વર્ષમાં બાર મહિના હોય છે.
\end{enumerate}
\sourcecorrection{OLMD-077}{મૂળમાં લેખકનું નામ ``Jaako Hintikka''
લખ્યું છે; ગ્રંથસૂચિ-વિવરણ મુજબ ``Jaakko Hintikka'' છે. અહીં નામ
સુધાર્યું છે.}
આધુનિક જ્ઞાનસંબંધી તર્કની શરૂઆત ઘણી વાર યાક્કો હિન્તિક્કાના
૧૯૬૨ના પુસ્તક~\emph{Knowledge and Belief}માં શોધાય છે. એ પુસ્તક
લખાયું ત્યારે તર્કમાં શક્ય જગતોનું અર્થવિજ્ઞાન વધુને વધુ વપરાવા
લાગ્યું હતું. હકીકતમાં, જ્ઞાનસંબંધી તર્કપદ્ધતિઓ બીજી મોડલ
તર્કપદ્ધતિઓનાં મોટા ભાગનાં એ જ અર્થવિજ્ઞાનલક્ષી સાધનો વાપરે છે,
પરંતુ તેમનું અર્થઘટન જુદી રીતે કરે છે. મુખ્ય ફેરફાર એ છે કે
\emph{સુલભતા સંબંધ} શું દર્શાવે છે એવું આપણે માનીએ છીએ.
જ્ઞાનસંબંધી તર્કમાં તે કોઈ પ્રકારની \emph{જ્ઞાનસંબંધી શક્યતા}
દર્શાવે છે. આપણે જોઈશું કે જે જ્ઞાનસંબંધી ખ્યાલનું નિદર્શન કરીએ
તેના પર સુલભતા સંબંધને મૂકવાની શરતો આધારિત હોય છે. અનુરૂપતા
સિદ્ધાંતને જ્ઞાનસંબંધી અર્થઘટન આપીએ ત્યારે શું થાય છે તે પણ
જોઈશું. ઉપરનાં ઉદાહરણોમાં બે કર્તાઓ, રિચાર્ડ અને ઓડ્રી, તથા
દરેક શું જાણે છે તેનો પરસ્પર સંબંધ આવે છે. આપણે વિચારીશું તે
જ્ઞાનસંબંધી તર્કપદ્ધતિઓ બહુ-કર્તા તર્કપદ્ધતિઓ છે, જેમાં આવી બાબતો
વ્યક્ત થઈ શકે છે. તેની સામે, એક-કર્તા જ્ઞાનસંબંધી તર્કપદ્ધતિ
માત્ર એક વ્યક્તિ શું જાણે છે કે માને છે તે વિશે વાત કરે છે.

\end{document}
